\scanpage{5}
\mspage{3}

\unsure{Est|est} équivalente à celle des W-algèbres locales, annulées par une puissance de l'idéal maximal de W, à extension résiduelle triviale. Ecrivant $W=\varprojlim W_n$ comme à l'accoutumé, dans le cas $p>0$, on trouve que la donnée d'un cristal de modules (absolu) sur k équivaut à celle d'un système projectif $(M_n)$ ``p-adique'' de modules $M_n$ sur les $W_n$; donc les cristaux de modules de présentation finie sur k forment une catégorie équivalente à celle des modules de type fini sur W. Si $p=0$, alors la catégorie des cristaux de modules sur k est équivalente à celle des \deleted{\illegible{1}} vectoriels sur $k=W$. Ces descriptions sont compatibles avec les notions de changement de base pour k variable. Elles se formulent évidemment pour toute autre espèce de cristaux, défini par une catégorie fibrée $\underline F$.

Dans le cas envisagé, le foncteur\deleted{\unsure{:}} ``valeur en S'', sur la catégorie des cristaux de présentation finie sur k, s'identifie au foncteur $\otimes_Wk$ sur la catégorie des modules de type fini sur W, foncteur qui (si $p>0$) n'est pas fidèle.

1.6. \uline{Exemples} 2: \inserted{S étale sur R.} Si $S=R$, la catégorie C admet R lui-même comme objet final, donc le foncteur ``valeur en R'' est une équivalence de la catégorie des cristaux sur R, de type $\underline F$ donné, avec la catégorie $\underline F_R$. En particulier, un cristal de modules sur $\operatorname{Spec}\underline Z$ est essentiellement la même chose qu'un $\underline Z$-module.

De façon un peu plus générale, si S est étale sur R, alors S est un objet final de C, et les cristaux de modules (disons) sur S, relativement à R, s'\unsure{identifient|ddentifient} aux Modules quasi-cohérents sur S.

1.7. \deleted{\uline{\unsure{Relations} avec la notion \unsure{d'objet|d'objets} \unsure{stratifié|stratifiés}}} \uline{Exemple} 3 : S \uline{un sous}-\deleted{\unsure{objet de|schéma}}\uline{préschéma de} R. Comme la notion de cristal relatif sur S ne change pas si on remplace R par un ouvert par lequel se factorise S, on peut supposer S fermé dans R, défini par un Idéal quasi-\unsure{cohhérent|cohérent} J. Soit $\deleted{R}\inserted{S}_n^{\deleted{\unsure{n}}}=V(J^{n+1})$ le n.ème voisinage infinitésimal de R dans S. Alors la famille des objets $\deleted{R}\inserted{S}_n$ de C est finale dans C, d'où on conclut facilement qu'un cristal de type $\underline F$ sur S s'identifie à une suite $(M_n)_{n\in\underline N}$ d'objets des $F_{\deleted{R}\inserted{S}_n}$ qui se recollent. En particulier, si R est localement noethérien, un cristal de modules sur S relativement à R, de présentation finie, s'identifie à un Module cohérent sur le complété formel de R le long de S.
\margin{4}
